% Research module: additive depth-baseline and amplitude entropy floor.
\section{A joint depth and amplitude cost below criticality}
\label{sec:joint-prefactor-entropy}
A polynomially weighted change of measure compares the entire subcritical variance chain with the canonical chain. It retains the depth cost while also retaining the sharp amplitude penalty.

\begin{lemma}[A variance-chain kernel domination]
\label{lem:joint-kernel-domination}
Let $N\ge2$, $0<\nu\le1$, and put
\[
 \beta=1-\nu,\qquad p=\frac{N\beta}{2\nu},\qquad
 \kappa(\nu)=1-\nu+\nu\log\nu,\qquad
 M_\nu=\exp\{-N\kappa(\nu)/(2\nu)\}.
\]
For the native scalar variance transition
\[
 q'=\frac1N\sum_{i=1}^N\tanh^2(\sqrt{\nu q}\,G_i)
\]
and every nonnegative measurable function $f$,
\[
 \E[(q')^p f(q')\mid q]
 \le M_\nu q^p\E[f(\widetilde q)\mid q],
 \qquad
 \widetilde q=\frac1N\sum_{i=1}^N\tanh^2(\sqrt q\,G_i).
\]
Thus the polynomially weighted subcritical transition, normalized by $M_\nu$, is dominated as a positive kernel by the canonical transition.
\end{lemma}
\begin{proof}
For $q>0$ let $S=N^{-1}\sum_iG_i^2$. Then $q'\le\nu qS$. Concavity of the logarithm, with tangent at $S=1/\nu$, gives
\[
 (\nu S)^p\le
 \exp\{p\nu(S-1/\nu)\}
 =\exp\left\{\frac\beta2\sum_iG_i^2-\frac{N\beta}{2\nu}\right\}.
\]
Multiplication by the standard Gaussian density changes its quadratic exponent to $-\nu\sum_iG_i^2/2$. The normalization factor is
\[
 \nu^{-N/2}e^{-N\beta/(2\nu)}
       =e^{-N\kappa(\nu)/(2\nu)}=M_\nu.
\]
Under the resulting product law, $G_i$ has variance $1/\nu$, so $\sqrt{\nu q}G_i$ becomes a canonical $N(0,q)$ preactivation. Multiplying by $q^p$ proves the inequality. The cases $q=0$ or $p=0$ follow directly, with the usual convention $q^0=1$.
\end{proof}

\begin{corollary}[A joint fractional and exponential moment]
\label{cor:joint-variance-moment}
For $3/2\le\theta<2$, let $k_\theta>0$ be the constant in the canonical heat supersolution. Then, for any native first layer with $q_1\le1$,
\[
 \E[q_H^p e^{\theta Nq_H}]
 \le\exp\left\{-\frac{N\kappa(\nu)(H-1)}{2\nu}
             +\frac{\theta N}{1+k_\theta(H-1)}\right\}.
\]
\end{corollary}
\begin{proof}
For $q>0$, divide the preceding lemma by $M_\nu q^p$ to obtain a subprobability transition kernel dominated by the canonical transition. Iteration preserves domination for every nonnegative test function. The polynomial weights telescope along a path, giving $M_\nu^{H-1}q_1^p$ times a canonical-chain expectation started at $q_1$. The canonical heat-supersolution estimate bounds its exponential moment uniformly over $q_1\le1$ by $\exp\{\theta N/[1+k_\theta(H-1)]\}$. Finally $q_1^p\le1$; integrate over the first layer.
\end{proof}

\begin{theorem}[An additive baseline-plus-amplitude entropy floor]
\label{thm:joint-prefactor-floor}
Let $N\ge2$, let $\mathcal X$ be any finite input set of size $m\ge1$, put $Z_*=\max_{x\in\mathcal X}|z_w(x)|$, and use the Gaussian tanh prior with $s=\nu/N$, $0<\nu\le1$. Fix $0<u<1$ and $0<t<2$. Choose $3/2\le\theta<2$ with $t^2<2\theta$, and write
\[
 L_{t,\theta}=\log\frac1{1-t^2/(2\theta)},\qquad
 B=\kappa(\nu)(H-1).
\]
Every comparison law $\rho$ satisfies
\begin{align*}
 2s\operatorname{KL}(\rho\|\mu)
 &\ge B\,\rho\{Z_*>u\}+2\nu t\,\E_\rho Z_* -\mathcal E,\\
 \mathcal E
 &=\frac{2\nu\theta}{1+k_\theta(H-1)}
   +\left(\beta+\frac\nu N\right)L_{t,\theta}
   +2\beta\log(1/u)
   +\frac{2\nu\log(2m)}N+2\nu t u.
\end{align*}
The expectation can be interpreted by monotone truncation. In particular, at $\beta=O(H^{-1/2})$ and fixed $u,t,\theta$, all error terms other than $2\nu t u$ vanish as $H,N\to\infty$ with $\log m=o(N)$.
\end{theorem}
\begin{proof}
At a fixed input the readout is conditionally $z=\sqrt{\nu q_H/N}\,G$, with an independent standard Gaussian $G$. Young's inequality gives
\[
 tN|z|\le\theta Nq_H+\frac{t^2\nu}{4\theta}G^2.
\]
Set $r=t^2\nu/(2\theta)<1$. Gaussian integration yields
\[
 \E|G|^{2p}e^{rG^2/2}
 =(1-r)^{-p-1/2}\E|G|^{2p}.
\]
Also $(\E|G|^{2p})^{1/(2p)}\le\sqrt{2p+2}$ for $p>0$ by the even-integer moment formula and monotonicity of moments. Since
\[
 (\nu/N)(2p+2)=\beta+2\nu/N\le1,
\]
we obtain
\[
 \E\bigl[|z|^{2p}e^{tN|z|}\bigr]
 \le(1-r)^{-p-1/2}\E[q_H^p e^{\theta Nq_H}].
\]
This inequality also holds when $p=0$ directly. On $|z|>u$, divide by $u^{2p}$. Corollary~\ref{cor:joint-variance-moment} and $1-r\ge1-t^2/(2\theta)$ then give
\[
 \E[e^{tN|z|}\mathbf1_{|z|>u}]
 \le\exp\left\{-\frac{NB}{2\nu}+\frac{N e}{2\nu}\right\},
\]
where
\[
 e=\frac{2\nu\theta}{1+k_\theta(H-1)}
    +\left(\beta+\frac\nu N\right)L_{t,\theta}
    +2\beta\log(1/u).
\]
The exponential of a maximum on its threshold event is at most the sum of the corresponding thresholded exponentials. Hence the same bound for $Z_*$ has $e$ replaced by $e+2\nu\log m/N$.

Let $A=\{Z_*>u\}$ and use the entropy test potential
\[
 V=\left(\frac{NB}{2\nu}+tNZ_*\right)\mathbf1_A.
\]
Its prior exponential moment is at most
$1+\exp\{N[e+2\nu\log m/N]/(2\nu)\}$, so its logarithm is at most $N[e+2\nu\log m/N]/(2\nu)+\log2$. The entropy variational inequality, multiplied by $2s=2\nu/N$, gives
\[
 2s\operatorname{KL}\ge B\rho(A)+2\nu t\E_\rho[Z_*\mathbf1_A]
                -e-2\nu\log(2m)/N.
\]
Finally $\E[Z_*\mathbf1_A]\ge\E Z_*-u$. This proves the theorem. For infinite moments, truncate the test potential and then pass to the limit.
\end{proof}
